\documentclass{csbeamer}

\usepackage{hyperref}
\usepackage{listings}
\usepackage{xcolor}
\usepackage{graphicx}
\usepackage{tikz}
\usetikzlibrary{positioning,shapes,arrows,calc}

% Python code styling
\lstset{
    language=Python,
    basicstyle=\ttfamily\small,
    keywordstyle=\color{blue}\bfseries,
    stringstyle=\color{red},
    commentstyle=\color{green!50!black}\itshape,
    showstringspaces=false,
    breaklines=true,
    frame=single,
    numbers=left,
    numberstyle=\tiny\color{gray}
}

% Course information
\title{Data Analysis and Visualization}
\course{CSCI 128: Coding for Problem Solving}
\courseshort{CSCI 128}
\term{Winter 2026}
\author{Dr. Jean-Alexis Delamer}

\begin{document}

\begin{frame}
\titlepage
\end{frame}

% ============================================================
\begin{frame}{Today's Topics}
    \begin{center}
        \Large From raw data to visual insight!
    \end{center}

    \vspace{1em}

    \textbf{Learning Objectives:}
    \begin{itemize}
        \item<1-> \textcolor{stfxblue}{Review:} loading and filtering CSV data
        \item<2-> \textcolor{marigold}{Group} and summarize data by category
        \item<3-> Compute \textcolor{stfxblue}{frequency counts}
        \item<4-> Visualize data with \textcolor{marigold}{matplotlib} bar charts and line plots
        \item<5-> \textcolor{stfxblue}{Compare groups} using statistics
        \item<6-> Build a complete \textcolor{marigold}{data analysis script}
        \item<7-> \textcolor{stfxblue}{Interpret results} and draw conclusions
    \end{itemize}

    \vspace{1em}

    \begin{center}
        \onslide<8->{\textit{Just like image filters transform pixels, charts \textcolor{marigold}{\textbf{transform numbers into stories!}}}}
    \end{center}
\end{frame}

% ============================================================
\begin{frame}{Review: Our Data Toolkit So Far}
    \begin{columns}
        \begin{column}{0.48\textwidth}
            \textbf{\textcolor{stfxblue}{Topic 7 --- What We Learned:}}
            \begin{itemize}
                \item<1-> Read CSV files with \texttt{csv.DictReader}
                \item<2-> Access rows and columns by name
                \item<3-> Compute \texttt{min}, \texttt{max}, and average
                \item<4-> Filter rows with \texttt{if} conditions
                \item<5-> Write results to a new CSV
            \end{itemize}
        \end{column}
        \begin{column}{0.48\textwidth}
            \onslide<6->{
            \textbf{\textcolor{marigold}{Topic 8 --- What We Add Today:}}
            \begin{itemize}
                \item<6-> Group rows by a \textcolor{marigold}{category}
                \item<7-> Count value \textcolor{stfxblue}{frequencies}
                \item<8-> Sort groups by computed values
                \item<9-> Create \textcolor{marigold}{bar charts} and \textcolor{stfxblue}{line plots}
                \item<10-> Produce charts as \texttt{.png} files
            \end{itemize}
            }
        \end{column}
    \end{columns}

    \vspace{1em}

    \onslide<11->{
    \begin{center}
        \textit{Think of Topic 7 as loading the image --- Topic 8 is applying the filter!}
    \end{center}
    }
\end{frame}

% ============================================================
\begin{frame}[fragile]{Asking Better Questions}
    \textbf{Imagine we have this weather dataset:}

    \vspace{0.5em}

    \begin{center}
    \begin{tabular}{|l|l|c|c|}
        \hline
        \textbf{City} & \textbf{Month} & \textbf{TempC} & \textbf{Rainfall} \\
        \hline
        Halifax   & January  & -5.2 & 110 \\
        Toronto   & January  & -3.7 & 62  \\
        Vancouver & January  &  4.1 & 154 \\
        Halifax   & July     & 19.8 & 94  \\
        Toronto   & July     & 22.3 & 74  \\
        Vancouver & July     & 18.0 & 36  \\
        \hline
    \end{tabular}
    \end{center}

    \vspace{0.8em}

    \textbf{Basic questions} (Topic 7):
    \begin{itemize}
        \item What is the highest temperature in the dataset?
    \end{itemize}

    \pause

    \vspace{0.4em}

    \textbf{\textcolor{marigold}{Better questions}} (Topic 8):
    \begin{itemize}
        \item<2-> Which \textcolor{stfxblue}{city} is warmest on average across all months?
        \item<3-> Which \textcolor{stfxblue}{month} has the most rainfall across all cities?
        \item<4-> How does temperature \textcolor{marigold}{change over the year} for each city?
    \end{itemize}

    \onslide<5->{\vspace{0.5em}
    \begin{center}
        \textit{Answering these requires \textcolor{marigold}{\textbf{grouping}} the data first.}
    \end{center}}
\end{frame}

% ============================================================
\begin{frame}{Grouping Data: The Concept}
    \textbf{Grouping = sorting rows into \textcolor{marigold}{buckets} by a shared value:}

    \vspace{0.8em}

    \begin{center}
    \begin{tikzpicture}[
        box/.style={draw, rounded corners, minimum width=2.2cm, minimum height=0.55cm, font=\small},
        bucket/.style={draw, fill=blue!10, rounded corners, minimum width=2.4cm, minimum height=2.2cm},
        label/.style={font=\small\bfseries},
        arrow/.style={->, thick, color=stfxblue}
    ]
        % Input rows
        \node[box, fill=red!10]   (r1) at (0, 3.2)  {Halifax, Jan, -5.2};
        \node[box, fill=green!10] (r2) at (0, 2.4)  {Toronto, Jan, -3.7};
        \node[box, fill=yellow!20](r3) at (0, 1.6)  {Vancouver, Jan, 4.1};
        \node[box, fill=red!10]   (r4) at (0, 0.8)  {Halifax, Jul, 19.8};
        \node[box, fill=green!10] (r5) at (0, 0.0)  {Toronto, Jul, 22.3};
        \node[box, fill=yellow!20](r6) at (0,-0.8)  {Vancouver, Jul, 18.0};

        % Group-by label
        \node[label] at (3.4, 3.6) {Group by \texttt{City}};

        % Buckets
        \node[bucket] (bH) at (3.2, 1.8) {};
        \node[label, above] at (3.2, 2.9) {\textcolor{red!70!black}{Halifax}};
        \node[font=\tiny] at (3.2, 2.3) {-5.2};
        \node[font=\tiny] at (3.2, 1.8) {19.8};
        \node[font=\tiny, color=gray] at (3.2, 1.3) {avg: 7.3};

        \node[bucket] (bT) at (5.8, 1.8) {};
        \node[label, above] at (5.8, 2.9) {\textcolor{green!50!black}{Toronto}};
        \node[font=\tiny] at (5.8, 2.3) {-3.7};
        \node[font=\tiny] at (5.8, 1.8) {22.3};
        \node[font=\tiny, color=gray] at (5.8, 1.3) {avg: 9.3};

        \node[bucket] (bV) at (8.4, 1.8) {};
        \node[label, above] at (8.4, 2.9) {\textcolor{orange!80!black}{Vancouver}};
        \node[font=\tiny] at (8.4, 2.3) {4.1};
        \node[font=\tiny] at (8.4, 1.8) {18.0};
        \node[font=\tiny, color=gray] at (8.4, 1.3) {avg: 11.1};

        % Arrows
        \draw[arrow, color=red!70!black]   (r1.east) to[out=0,in=180] (bH.west);
        \draw[arrow, color=green!50!black] (r2.east) to[out=0,in=180] (bT.west);
        \draw[arrow, color=orange!80!black](r3.east) to[out=0,in=180] (bV.west);
        \draw[arrow, color=red!70!black]   (r4.east) to[out=0,in=180] (bH.west);
        \draw[arrow, color=green!50!black] (r5.east) to[out=0,in=180] (bT.west);
        \draw[arrow, color=orange!80!black](r6.east) to[out=0,in=180] (bV.west);
    \end{tikzpicture}
    \end{center}

    \vspace{0.3em}
    \begin{center}
        Each bucket holds all rows that share the same \textcolor{stfxblue}{\textbf{City}} value.\\
        Then we can compute stats \textbf{per group}.
    \end{center}
\end{frame}

% ============================================================
\begin{frame}[fragile]{Grouping with a Dictionary}
    \textbf{The grouping pattern in Python:}

    \begin{lstlisting}
import csv

# Load the data
with open('weather.csv', newline='') as f:
    reader = csv.DictReader(f)
    data = list(reader)

# Step 1: Build groups (dict of lists)
groups = {}
for row in data:
    key = row['City']
    if key not in groups:
        groups[key] = []          # start a new bucket
    groups[key].append(float(row['TempC']))

# Step 2: Compute per-group averages
averages = {}
for city, temps in groups.items():
    averages[city] = sum(temps) / len(temps)

# Step 3: Display results
for city, avg in averages.items():
    print(f"{city}: {avg:.1f} C average")
    \end{lstlisting}

    \pause

    \begin{center}
        \textit{The \textcolor{marigold}{\texttt{if key not in groups}} check creates the bucket \textbf{on first encounter}.}
    \end{center}
\end{frame}

% ============================================================
\begin{frame}{Try It: Group by Category}
    \begin{center}
        \Large \textbf{Activity Time!}
    \end{center}

    \vspace{1em}

    \textbf{You have a file \texttt{sports.csv} with columns:}\\
    \texttt{Team, Player, Goals, Assists, Wins}

    \vspace{0.8em}

    \textbf{Tasks:}
    \begin{enumerate}
        \item Load the CSV with \texttt{csv.DictReader}
        \item Group the rows by \texttt{Team}
        \item For each team, sum the \texttt{Wins} column
        \item Print the team name and total wins for each team
        \item Print which team has the \textbf{most} wins overall
    \end{enumerate}

    \vspace{0.8em}

    \textbf{Bonus:} Also compute the average goals per team and rank teams by it.
\end{frame}

% ============================================================
\begin{frame}[fragile]{Frequency Counts}
    \textbf{How many times does each value appear?}

    \vspace{0.5em}

    \begin{columns}
        \begin{column}{0.55\textwidth}
    \begin{lstlisting}
import csv

with open('grades.csv', newline='') as f:
    reader = csv.DictReader(f)
    data = list(reader)

# Count occurrences of each grade
counts = {}
for row in data:
    grade = row['LetterGrade']
    if grade not in counts:
        counts[grade] = 0
    counts[grade] += 1

# Print the frequency table
for grade, n in sorted(counts.items()):
    bar = '#' * n
    print(f"{grade}: {bar} ({n})")
    \end{lstlisting}
        \end{column}
        \begin{column}{0.42\textwidth}
            \onslide<2->{
            \textbf{Sample output:}
            \begin{center}
            \begin{tabular}{ll}
                \texttt{A:} & \texttt{\#\#\#\#\#\#\#\#\#\#\#\# (12)} \\
                \texttt{B:} & \texttt{\#\#\#\#\#\#\#\#\#\#\#\#\#\#\#\#\#\# (18)} \\
                \texttt{C:} & \texttt{\#\#\#\#\#\#\#\#\# (9)} \\
                \texttt{D:} & \texttt{\#\#\#\#\# (5)} \\
                \texttt{F:} & \texttt{\#\# (2)} \\
            \end{tabular}
            \end{center}

            \vspace{0.5em}
            \textit{This is the data behind a \textcolor{marigold}{\textbf{bar chart}} --- we'll plot it soon!}
            }
        \end{column}
    \end{columns}
\end{frame}

% ============================================================
\begin{frame}[fragile]{Sorting Groups by Value}
    \textbf{Ranking your groups from highest to lowest:}

    \begin{lstlisting}
# averages is a dict: {'Halifax': 7.3, 'Toronto': 9.3, ...}

# Sort by value, descending
ranked = sorted(averages.items(), key=lambda x: x[1], reverse=True)

print("Cities ranked by average temperature:")
for rank, (city, avg) in enumerate(ranked, start=1):
    print(f"  {rank}. {city}: {avg:.1f} C")
    \end{lstlisting}

    \pause

    \textbf{How \texttt{sorted()} works here:}
    \begin{itemize}
        \item<2-> \texttt{averages.items()} gives pairs like \texttt{('Halifax', 7.3)}
        \item<3-> \texttt{key=lambda x: x[1]} sorts by the \textcolor{marigold}{second element} (the value)
        \item<4-> \texttt{reverse=True} orders from \textcolor{stfxblue}{largest to smallest}
        \item<5-> \texttt{enumerate(..., start=1)} adds a \textcolor{marigold}{1-based rank counter}
    \end{itemize}

    \pause

    \begin{center}
        \textit{This same pattern works for \textbf{any} dictionary of numeric values.}
    \end{center}
\end{frame}

% ============================================================
\begin{frame}{Introducing Matplotlib}
    \textbf{What is \textcolor{marigold}{matplotlib}?}

    \vspace{0.5em}

    \begin{itemize}
        \item<1-> The most popular Python library for creating \textcolor{stfxblue}{charts and graphs}
        \item<2-> Produces publication-quality figures as PNG, PDF, SVG, and more
        \item<3-> Used by scientists, journalists, and data analysts worldwide
    \end{itemize}

    \vspace{0.8em}

    \onslide<4->{\textbf{Installing it} (it is not in the standard library):}
    \onslide<4->{\begin{center}\texttt{pip install matplotlib}\end{center}}

    \vspace{0.5em}

    \onslide<5->{\textbf{Importing it} (the standard alias):}
    \onslide<5->{\begin{center}\texttt{import matplotlib.pyplot as plt}\end{center}}

    \vspace{0.8em}

    \onslide<6->{\textbf{The basic pattern:}
    \begin{enumerate}
        \item<6-> Prepare your data as Python \textcolor{marigold}{lists}
        \item<7-> Call a chart function: \texttt{plt.bar()}, \texttt{plt.plot()}, \texttt{plt.hist()}
        \item<8-> Add labels with \texttt{plt.title()}, \texttt{plt.xlabel()}, \texttt{plt.ylabel()}
        \item<9-> Save with \texttt{plt.savefig('chart.png')} and/or display with \texttt{plt.show()}
    \end{enumerate}}
\end{frame}

% ============================================================
\begin{frame}[fragile]{Your First Bar Chart}
    \textbf{Plotting grade counts as a bar chart:}

    \begin{lstlisting}
import matplotlib.pyplot as plt

# Data (from our frequency count earlier)
labels = ['A', 'B', 'C', 'D', 'F']
counts = [12, 18, 9, 5, 2]

# Create the bar chart
plt.bar(labels, counts)

# Add labels and a title
plt.title('Grade Distribution')
plt.xlabel('Grade')
plt.ylabel('Number of Students')

# Save to a file, then display
plt.savefig('grades.png')
plt.show()
    \end{lstlisting}

    \pause

    \begin{center}
        \textit{Just six meaningful lines of code produce a complete, shareable chart!}
    \end{center}
\end{frame}

% ============================================================
\begin{frame}[fragile]{Customizing Charts}
    \textbf{Making charts clearer and more professional:}

    \begin{lstlisting}
import matplotlib.pyplot as plt

labels = ['A', 'B', 'C', 'D', 'F']
counts = [12, 18, 9, 5, 2]

# Set figure size (width, height in inches)
plt.figure(figsize=(8, 5))

# Bar color and edge color
plt.bar(labels, counts,
        color='steelblue',
        edgecolor='black')

# Labels and title
plt.title('Grade Distribution --- CSCI 128', fontsize=14)
plt.xlabel('Letter Grade', fontsize=12)
plt.ylabel('Number of Students', fontsize=12)

# Add a subtle horizontal grid (easier to read bar heights)
plt.grid(axis='y', alpha=0.3)

plt.tight_layout()          # prevent clipping
plt.savefig('grades.png', dpi=150)
plt.show()
    \end{lstlisting}
\end{frame}

% ============================================================
\begin{frame}[fragile]{Line Plots: Showing Change Over Time}
    \textbf{A line plot is ideal when data has \textcolor{marigold}{order} (months, years, steps):}

    \begin{lstlisting}
import matplotlib.pyplot as plt

months = ['Jan', 'Feb', 'Mar', 'Apr', 'May', 'Jun',
          'Jul', 'Aug', 'Sep', 'Oct', 'Nov', 'Dec']
temps  = [-5.2, -4.1, 0.3, 5.8, 12.1, 17.4,
           19.8, 19.2, 14.6, 8.3, 2.5, -2.9]

plt.figure(figsize=(10, 5))

# plot() draws a line; marker='o' adds a dot at each point
plt.plot(months, temps, marker='o',
         color='steelblue', linewidth=2, markersize=6)

plt.title('Average Monthly Temperature --- Halifax')
plt.xlabel('Month')
plt.ylabel('Temperature (C)')
plt.grid(axis='y', alpha=0.3)
plt.tight_layout()
plt.savefig('temperature_halifax.png', dpi=150)
plt.show()
    \end{lstlisting}
\end{frame}

% ============================================================
\begin{frame}[fragile]{Connecting CSV Data to a Chart}
    \textbf{The full pipeline: CSV $\rightarrow$ lists $\rightarrow$ chart:}

    \begin{lstlisting}
import csv
import matplotlib.pyplot as plt

# Step 1: Load CSV
with open('city_avg_temps.csv', newline='') as f:
    reader = csv.DictReader(f)
    data = list(reader)

# Step 2: Extract two columns into parallel lists
cities = []
avg_temps = []
for row in data:
    cities.append(row['City'])
    avg_temps.append(float(row['AvgTempC']))

# Step 3: Pass lists directly to plt.bar()
plt.figure(figsize=(8, 5))
plt.bar(cities, avg_temps, color='coral', edgecolor='black')
plt.title('Average Annual Temperature by City')
plt.xlabel('City')
plt.ylabel('Average Temperature (C)')
plt.grid(axis='y', alpha=0.3)
plt.tight_layout()
plt.savefig('city_temps.png', dpi=150)
plt.show()
    \end{lstlisting}

    \pause
    \begin{center}
        \textit{The \textcolor{marigold}{\textbf{same two lists}} that fed \texttt{print()} now feed \texttt{plt.bar()} --- no other change!}
    \end{center}
\end{frame}

% ============================================================
\begin{frame}{Try It: Visualize Your Dataset}
    \begin{center}
        \Large \textbf{Activity Time!}
    \end{center}

    \vspace{1em}

    \textbf{Choose any CSV dataset you have} (sports, weather, or grades).

    \vspace{0.8em}

    \textbf{Tasks:}
    \begin{enumerate}
        \item Load the CSV with \texttt{csv.DictReader}
        \item Choose a \textbf{category column} to group by (e.g. Team, City, Grade)
        \item Compute a \textbf{numeric summary} per group (total, average, or count)
        \item Build two parallel lists: one for labels, one for values
        \item Create a bar chart with a title, x-label, and y-label
        \item Save the chart as \texttt{my\_chart.png}
    \end{enumerate}

    \vspace{0.8em}

    \textbf{Bonus:} Customize the bar color and add a grid. Try sorting the bars from tallest to shortest before plotting.
\end{frame}

% ============================================================
\begin{frame}[fragile]{Comparing Two Groups}
    \textbf{Side-by-side bars let us \textcolor{marigold}{compare two series} at once:}

    \begin{lstlisting}
import matplotlib.pyplot as plt
import numpy as np

cities = ['Halifax', 'Toronto', 'Vancouver']
summer_temps = [19.8, 22.3, 18.0]
winter_temps = [-5.2, -3.7,  4.1]

x = np.arange(len(cities))   # [0, 1, 2]
width = 0.35                  # width of each bar

plt.figure(figsize=(8, 5))
plt.bar(x - width/2, summer_temps, width,
        label='July', color='coral',      edgecolor='black')
plt.bar(x + width/2, winter_temps, width,
        label='January', color='steelblue', edgecolor='black')

plt.title('Summer vs. Winter Temperatures')
plt.xlabel('City')
plt.ylabel('Temperature (C)')
plt.xticks(x, cities)         # replace 0,1,2 with city names
plt.legend()
plt.grid(axis='y', alpha=0.3)
plt.tight_layout()
plt.savefig('comparison.png', dpi=150)
plt.show()
    \end{lstlisting}
\end{frame}

% ============================================================
\begin{frame}[fragile]{Histograms: Distribution of a Single Column}
    \textbf{A histogram shows how values are \textcolor{marigold}{spread out} across a range:}

    \begin{lstlisting}
import csv
import matplotlib.pyplot as plt

# Load exam scores
with open('scores.csv', newline='') as f:
    reader = csv.DictReader(f)
    scores = [float(row['Score']) for row in reader]

plt.figure(figsize=(8, 5))

# bins=10 divides the range into 10 equal intervals
plt.hist(scores, bins=10, color='steelblue',
         edgecolor='black')

plt.title('Distribution of Exam Scores')
plt.xlabel('Score')
plt.ylabel('Number of Students')
plt.grid(axis='y', alpha=0.3)
plt.tight_layout()
plt.savefig('score_histogram.png', dpi=150)
plt.show()
    \end{lstlisting}

    \pause

    \textbf{Sound familiar?} In Topic 3 we asked ``how many pixels have this brightness?'' ---
    a histogram answers \textcolor{marigold}{\textit{exactly}} that question, for \textit{any} column of numbers.
\end{frame}

% ============================================================
\begin{frame}[fragile]{A Complete Analysis Script}
    \textbf{Putting it all together in one cohesive program:}

    \begin{lstlisting}
import csv
import matplotlib.pyplot as plt

# 1. Load CSV
with open('weather.csv', newline='') as f:
    data = list(csv.DictReader(f))

# 2. Group by City, collect temperatures
groups = {}
for row in data:
    city = row['City']
    if city not in groups:
        groups[city] = []
    groups[city].append(float(row['TempC']))

# 3. Compute average per group
averages = {city: sum(t)/len(t) for city, t in groups.items()}

# 4. Sort from warmest to coldest
ranked = sorted(averages.items(), key=lambda x: x[1], reverse=True)
labels = [item[0] for item in ranked]
values = [item[1] for item in ranked]

# 5. Save a bar chart
plt.figure(figsize=(9, 5))
plt.bar(labels, values, color='coral', edgecolor='black')
plt.title('Average Annual Temperature by City')
plt.xlabel('City')
plt.ylabel('Average Temperature (C)')
plt.grid(axis='y', alpha=0.3)
plt.tight_layout()
plt.savefig('avg_temps.png', dpi=150)
print("Chart saved to avg_temps.png")
    \end{lstlisting}
\end{frame}

% ============================================================
\begin{frame}{Interpreting Results}
    \textbf{Writing code is only half the job --- reading results carefully matters too.}

    \vspace{0.8em}

    \begin{itemize}
        \item<1-> \textcolor{stfxblue}{\textbf{Correlation $\neq$ causation}} \\
            {\small Two things moving together does not mean one \textit{causes} the other.}

        \vspace{0.4em}

        \item<2-> \textcolor{marigold}{\textbf{Watch for outliers}} \\
            {\small One extreme value can distort an average dramatically. Always check \texttt{min} and \texttt{max} alongside the mean.}

        \vspace{0.4em}

        \item<3-> \textcolor{stfxblue}{\textbf{Sample size matters}} \\
            {\small An average based on 3 rows is far less reliable than one based on 3\,000 rows. Report \textit{n} alongside your statistics.}

        \vspace{0.4em}

        \item<4-> \textcolor{marigold}{\textbf{What story does the data tell?}} \\
            {\small A good analysis ends with a plain-language sentence: \textit{``Vancouver has the highest average annual temperature of the three cities at 11.1\,\textdegree{}C.''}}

        \vspace{0.4em}

        \item<5-> \textcolor{stfxblue}{\textbf{Be honest about limitations}} \\
            {\small Is the dataset complete? Is it representative? Was data collected consistently?}
    \end{itemize}
\end{frame}

% ============================================================
\begin{frame}{Quick Self-Check}
    \textbf{Can you answer these?}

    \vspace{1em}

    \begin{enumerate}
        \item What Python data structure do we use to store per-group lists, and why?
        \item What does \texttt{key=lambda x: x[1]} mean in \texttt{sorted()}?
        \item Name two differences between \texttt{plt.bar()} and \texttt{plt.plot()}.
        \item When is a histogram more useful than a bar chart?
        \item You compute that Team A has a higher average score than Team B. Can you conclude Team A is better? What else would you need to know?
    \end{enumerate}

    \vspace{1em}

    \pause

    \begin{center}
        \textit{If any of these feel uncertain, review those slides or ask questions!}
    \end{center}
\end{frame}

% ============================================================
\begin{frame}{Real-World Applications}
    \textbf{The skills from today power real careers and projects:}

    \vspace{0.8em}

    \begin{itemize}
        \item<1-> \textbf{Sports analytics:} Grouping player stats by team, season, or position to find patterns and inform trades
        \item<2-> \textbf{Climate science:} Aggregating temperature and rainfall data by region and decade to study long-term trends
        \item<3-> \textbf{Public health:} Comparing disease rates across age groups or geographic areas to prioritize interventions
        \item<4-> \textbf{Business intelligence:} Breaking down sales figures by product category, region, or quarter for executive dashboards
        \item<5-> \textbf{Data journalism:} Turning government CSV datasets into charts that tell stories for a general audience
        \item<6-> \textbf{Machine learning pipelines:} Exploratory data analysis (EDA) is the mandatory first step before training any model
    \end{itemize}

    \vspace{0.8em}

    \onslide<7->{
    \begin{center}
        \textit{Every one of these domains uses the \textcolor{marigold}{\textbf{same group-summarize-visualize}} loop you practised today.}
    \end{center}
    }
\end{frame}

% ============================================================
\begin{frame}{For Next Class}
    \textbf{Before next session:}
    \begin{itemize}
        \item Complete the data analysis lab: load a CSV, group it, and produce at least two different charts
        \item Make sure \texttt{matplotlib} is installed in your Python environment
        \item Experiment: try changing \texttt{bins} in \texttt{plt.hist()} and observe what happens
        \item Think about: what would happen if your CSV had \textit{thousands} of rows? What would change?
    \end{itemize}

    \vspace{1em}

    \textbf{Next session --- Topic 9: Bigger Programs}
    \begin{itemize}
        \item Writing and calling your own \textcolor{marigold}{functions}
        \item Organizing code into \textcolor{stfxblue}{modules}
        \item Parameters, return values, and scope
        \item Structuring a multi-file project
        \item Reusing your analysis code across datasets
    \end{itemize}

    \vspace{0.8em}

    \begin{center}
        \textbf{Questions?}
    \end{center}
\end{frame}

\end{document}
